\section{Problema educacional e justificativa}

No ensino de Computação voltado ao Ensino Médio, o problema educacional investigado nesta pesquisa está na distância entre reconhecer elementos visuais de uma ferramenta digital e compreender a lógica que produz seus resultados. Embora a compreensão de um modelo de \textit{machine learning} dependa do contato com dados e de noções básicas de estatística e probabilidade, o estudante pode identificar nomes técnicos, gráficos e controles na tela sem conseguir explicar como dados, finalidade da tarefa, regras e parâmetros se relacionam com o resultado.

Já existem iniciativas voltadas à Educação Básica. Ao analisarem 30 unidades didáticas dedicadas ao ensino de \textit{machine learning} na escola, \citeonline{marques2020teaching} observaram que várias delas cobrem apenas processos mais acessíveis, como a gestão de dados, ou apresentam o aprendizado e o teste dos modelos de forma abstrata. Nessas unidades, parte do funcionamento interno dos modelos permanece oculta. Enquanto esse contato restrito à gestão deixa de mostrar como os dados e os parâmetros alteram o resultado do modelo, a revisão de \citeonline{martins2023findings} no Ensino Médio indica que os estudantes conseguem compreender conceitos básicos da área e aplicá-los em tarefas com algoritmos. Essa revisão aponta a necessidade de avaliações mais rigorosas. Diante desse cenário investigado na literatura recente da área, o trabalho pedagógico na escola pede atividades que vinculem os dados e a finalidade da tarefa ao modo como regras e parâmetros modificam o resultado.

Quando os dados chegam prontos às aulas e as relações entre eles não são trabalhadas, o estudante tende a aprender os conceitos de \textit{machine learning} apenas pelo nome, sem conseguir passar do reconhecimento do conceito ao seu uso na explicação de uma situação. O vocabulário básico acaba repetido sem conexão real. Sem saber explicar o que muda quando um parâmetro é alterado durante a atividade escolar, um estudante pode dizer que um modelo foi treinado e reconhecer uma reta no gráfico sem relacioná-la a uma previsão.

Cultura digital e Computação complementam-se na formação escolar. Enquanto a discussão cultural é indispensável para problematizar impactos sociais que abrangem desde vieses e riscos até questões de autoria, o recorte no ensino de Computação exige que o estudante também investigue na prática uma representação de dados e uma regra computacional. Sem reduzir questões éticas a uma equação ou a um gráfico no trabalho escolar, considera-se que a discussão sobre limites passa a ter outra base quando o estudante reconhece que modelos resultam de decisões humanas de modelagem.

Esse processo começa na própria origem dos dados. Formulando previsões antes de alterar dados ou parâmetros para observar e explicar o efeito produzido na atividade, o estudante que monta uma tabela com medidas da própria turma passa a saber de onde vêm os números e o que cada coluna representa. Essa experiência torna a leitura do gráfico menos mecânica. Quando vê em seguida uma reta se deslocar após a mudança nos dados ou uma fronteira aproximar-se de uma classe, o estudante encontra evidências visíveis para relacionar na explicação o efeito observado à alteração realizada, aproximando-se das abordagens práticas baseadas em problemas ou projetos relatadas com resultados positivos por \citeonline{martins2023findings}. Compreendendo de onde vêm os dados trabalhados e o que muda com o ajuste de um parâmetro, o estudante encontra elementos para questionar o resultado que recebe em vez de apenas aceitá-lo.

O produto educacional reúne sequência didática e plataforma. Para responder a esse problema educacional no Ensino Médio sem exigir a mesma formalização matemática em cada etapa, a plataforma em português brasileiro organiza cinco módulos introdutórios como conteúdos centrais: regressão linear, máquina de vetores de suporte (\textit{SVM}), \textit{k}-vizinhos mais próximos (\textit{k}-NN), árvore de decisão e \textit{k}-means. Percorrendo esses módulos para examinar na prática a relação entre dados e resultados, a turma trabalha desde previsões numéricas com linha de tendência e medição de erro até a separação de classes por fronteira, distância ou regras e o agrupamento de dados não rotulados.

Usar a plataforma não comprova a aprendizagem. Como o potencial das ferramentas visuais depende de como as representações são desenhadas e das tarefas cognitivas pedidas aos estudantes \cite{ainsworth2006deft}, cada módulo partirá de uma situação-problema e da origem dos dados para que a turma levante uma hipótese antes da manipulação controlada, avançando da previsão inicial ao experimento prático e à explicação do que foi observado. Serão oito encontros de cerca de 50 minutos. Com os resultados do experimento retomados pela mediação docente nos moldes da avaliação formativa discutida por \citeonline{black1998assessment}, a turma poderá organizar no segundo encontro uma tabela simples antes de estudar regressão linear e trabalhará nos demais módulos com cenários sintéticos fornecidos pela plataforma. A mudança na margem exige discussão em aula. Caso o estudante que aproxima a fronteira de uma classe na tela deixe de discutir por que a margem mudou, a manipulação do controle visual pode ficar desvinculada do conceito que se pretende ensinar.

Diante da limitação de tempo das aulas e da variação de infraestrutura entre as escolas da Educação Básica, o trabalho docente em turmas heterogêneas demanda recursos pedagógicos cujas condições de aplicação estejam claramente documentadas. Uma interface isolada não resolve dificuldades de aprendizagem. No contexto de um programa profissional como o ProfComp, a justificativa do estudo está em vincular esse problema escolar ao conteúdo de Computação por meio de um produto utilizável, cuja análise considera seu uso em aula sob a mediação direta do professor. Tornando visíveis algumas relações envolvidas no funcionamento dos modelos sem exigir nas aulas o aprofundamento próprio dos cursos técnicos, a proposta busca aproximar o ensino introdutório de \textit{machine learning} da realidade do Ensino Médio regular. A contribuição esperada dependerá das condições de aplicação. Sem antecipar generalizações que esta pesquisa não poderá sustentar, a documentação das condições locais poderá ajudar outros docentes a avaliar se uma proposta semelhante serve ao seu próprio contexto escolar.

\section{Pergunta de pesquisa}

Que mudanças e permanências podem ser identificadas nas explicações de estudantes do Ensino Médio sobre fundamentos introdutórios de \textit{machine learning} ao longo de uma sequência didática apoiada por uma plataforma interativa?

A análise acompanha diferentes etapas da sequência. Dirigindo a atenção para o modo como os estudantes justificam escolhas durante as atividades, a pergunta busca examinar como eles interpretam gráficos ou tabelas ao formular previsões e analisar os resultados observados em aula. Em vez de avaliar o efeito isolado da plataforma digital sobre o desempenho da turma, o estudo examinará o que muda ou permanece nas dificuldades de relacionar os dados e a finalidade da tarefa à alteração de regras ou parâmetros e aos resultados produzidos pelo modelo. Como os registros discentes serão elaborados em momentos definidos ao longo dos oito encontros previstos para a intervenção escolar, a investigação considerará produções construídas na combinação entre o uso individual da plataforma, a mediação docente e a discussão coletiva. Esse percurso respeitará as condições reais da turma.
